% =============================================================================
\section{Implementação}
\label{sec:implementacao}

% -----------------------------------------------------------------------------
\subsection{Representação dos Indivíduos (Genótipo)}
\label{sec:impl_representacao}

Os indivíduos são estruturados como árvores sintáticas recursivas através da
classe \code{Node(value, children)}, encapsuladas na entidade
\code{Individual}. O genótipo é composto por:

\begin{itemize}
  \item \textbf{Nós Terminais ($\mathcal{T}$):} constituídos pelas variáveis
        do problema ($x_0, x_1, \dots, x_{d-1}$) e por Constantes Efêmeras
        Aleatórias (ERC). Quando uma folha é sorteada, uma constante é
        gerada com probabilidade $p_{\text{const}} = 0.3$, amostrada
        uniformemente no intervalo $[-1.0, 1.0]$.
  \item \textbf{Nós Internos / Funções ($\mathcal{F}$):} aplicam operações
        aritméticas ou não-lineares aos nós filhos.
  \item \textbf{Forma Canônica:} cada indivíduo gera uma assinatura
        algébrica canônica normalizada (\code{canonical\_}\allowbreak\code{repr()}), que comuta operandos
        em nós simétricos (ex.: $x_0 + x_1 \equiv x_1 + x_0$), viabilizando a
        detecção exata de clones sintáticos na população.
\end{itemize}

\paragraph{Critério de Limite de Tamanho.}
O enunciado especifica o ``tamanho máximo do indivíduo igual a~7''.
Adotou-se \textbf{profundidade máxima 7} (na qual uma folha isolada possui
profundidade~1 e a raiz com filhos tem profundidade $1 + \max(d_i)$). Em
árvores estritamente binárias, esse teto comporta até $2^7 - 1 = 127$ nós.
Qualquer operador que gere descendentes com profundidade $>7$ tem a operação
rejeitada e reamostrada (até 20 tentativas); caso o limite persista em ser
violado, o pai original é preservado por reprodução.

% -----------------------------------------------------------------------------
\subsection{Função de Avaliação (Fitness) e Desacoplamento Numérico}
\label{sec:impl_fitness}

A aptidão básica a ser minimizada é a Raiz do Erro Quadrático Médio (RMSE):
\begin{equation}
  f(\mathit{Ind}) = \sqrt{\frac{1}{N}\sum_{i=1}^{N}\big(\mathrm{Eval}(\mathit{Ind},\mathbf{x}_i)-y_i\big)^2 },
  \label{eq:rmse}
\end{equation}
em que $\mathbf{x}_i$ é a $i$-ésima linha da matriz $X \in \mathbb{R}^{N \times d}$
e $y_i$ é o valor alvo correspondente. A avaliação é totalmente vetorizada: a
árvore executa diretamente sobre arrays NumPy em uma única passada. Quaisquer
ocorrências de indefinições numéricas (\code{NaN} ou \code{Inf}) são
interceptadas e penalizadas com aptidão de $10^{12}$.

\paragraph{Escalonamento Linear Ótimo (\emph{Linear Scaling}).}
Em problemas numéricos contínuos, forçar a árvore a buscar simultaneamente a
forma funcional $g(X)$ e a escala/deslocamento dos dados cria barreiras severas
na paisagem de busca. Para contornar essa dificuldade, implementou-se o
\emph{Linear Scaling}~\cite{keijzer2003improving}. A resposta ajustada é dada
por $\hat{y} = \alpha + \beta g(X)$, em que os coeficientes ótimos são
resolvidos analiticamente por Mínimos Quadrados Ordinários (OLS) em tempo
$\mathcal{O}(N)$ sobre o conjunto de treino:
\begin{equation}
  \beta = \frac{\sum_{i=1}^{N} (y_i - \bar{y})(g_i - \bar{g})}{\sum_{i=1}^{N} (g_i - \bar{g})^2},
  \qquad \alpha = \bar{y} - \beta\,\bar{g}.
  \label{eq:ls}
\end{equation}
Caso a variância de $g(X)$ seja nula ($\sum (g_i-\bar{g})^2 < 10^{-12}$), fixa-se
$\beta = 0$ e $\alpha = \bar{y}$. No conjunto de teste, a expressão é avaliada
utilizando exatamente os valores de $\alpha$ e $\beta$ congelados durante o
treino.

% -----------------------------------------------------------------------------
\subsection{Conjunto de Primitivas e Tratamento de Descontinuidades}
\label{sec:impl_primitivas}

O conjunto padrão da Fase~1 utiliza $\mathcal{F}_{\text{canônico}} = \{ +, -, \times, \div \}$,
com a \textbf{divisão protegida de Koza}~\cite{koza1992genetic}:
\begin{equation}
  \mathrm{div}(a,b) =
  \begin{cases}
    1.0, & |b| < 10^{-6},\\
    a/b, & \text{caso contrário.}
  \end{cases}
  \label{eq:div}
\end{equation}
Conforme demonstrado adiante nos experimentos, essa formulação clássica
introduz falésias artificiais quando $10^{-6} < |b| \ll 1$, gerando polos
assintóticos fora da amostra de treino. Para solucionar essa deficiência na
Fase~2, implementou-se o \textbf{\emph{Analytic Quotient} (AQ)}~\cite{ni2013analytic}:
\begin{equation}
  \mathrm{aq}(a,b) = \frac{a}{\sqrt{1 + b^2}},
  \label{eq:aq}
\end{equation}
que é infinitamente diferenciável e contínuo em todo o $\mathbb{R}^2$,
eliminando assíntotas verticais. Para problemas com dinâmicas não-lineares
complexas e respostas amortecidas, foram também disponibilizadas primitivas
unárias protegidas: $\log(|a| + 10^{-6})$, $\sqrt{|a|}$ e $\exp(a)$ (com corte
em $[-50, 50]$ para prevenir \emph{overflow}).

% -----------------------------------------------------------------------------
\subsection{População Inicial}
\label{sec:impl_init}

Foram implementados três métodos clássicos de geração:
\begin{enumerate}
  \item \textbf{\emph{Full}:} nós internos são sorteados até a profundidade
        alvo $d$, gerando árvores perfeitamente simétricas com folhas na
        profundidade máxima.
  \item \textbf{\emph{Grow}:} nós internos e folhas são sorteados com igual
        probabilidade desde a raiz, gerando topologias assimétricas e de
        tamanhos variados.
  \item \textbf{\emph{Ramped Half-and-Half} (Padrão adotado):} a população é
        particionada equitativamente entre as profundidades $d \in \{1, \dots, \text{max\_size}\}$.
        Em cada faixa, $50\%$ dos indivíduos são gerados via \emph{full} e
        $50\%$ via \emph{grow}. Esse mecanismo produz árvores com tamanho
        médio inicial de $26.2$ nós e profusão de material genético (cerca
        de $89.5\%$ a $95.1\%$ de indivíduos únicos na geração~0).
\end{enumerate}

% -----------------------------------------------------------------------------
\subsection{Operadores Genéticos de Variação}
\label{sec:impl_operadores}

A cada geração, o processo de reprodução opera em pares sobre cópias imutáveis
dos progenitores:
\begin{itemize}
  \item \textbf{\emph{Crossover} de Subárvore:} seleciona um ponto de corte em
        cada progenitor e permuta as subárvores correspondentes. Para evitar
        trocas triviais entre folhas, adota-se uma probabilidade de $90\%$ de
        sortear nós internos e $10\%$ para terminais.
  \item \textbf{Mutação de Subárvore:} sorteia um nó aleatório da árvore e o
        substitui por uma nova subárvore gerada pelo algoritmo \emph{grow}.
  \item \textbf{Mutação Pontual:} substitui uma função por outra de mesma
        aridade ou altera uma variável/constante, mantendo intacta a estrutura.
  \item \textbf{Mutação Mista (\emph{Mixed}):} introduzida na Fase~2 para
        suavizar a paisagem epistática, sorteia entre três modalidades:
        (i)~$35\%$ de perturbação gaussiana direta sobre constantes ERC
        ($c \leftarrow c + \mathcal{N}(0, \sigma^2)$, com $\sigma = 0.25\max(1, |c|)$);
        (ii)~$35\%$ de mutação pontual; e (iii)~$30\%$ de micro-subárvore
        \emph{grow}.
\end{itemize}
No fluxo geracional padrão, com probabilidade $\pc$ aplica-se o \emph{crossover}
(ou clonagem dos pais caso contrário) e, em seguida, cada descendente gerado
está sujeito à mutação com probabilidade $\pmut$.

% -----------------------------------------------------------------------------
\subsection{Mecanismos de Seleção, Elitismo e Controle de Bloat}
\label{sec:impl_selecao}

O arcabouço disponibiliza três paradigmas de seleção:

\paragraph{1. Seleção por Roleta (Proporcional à Aptidão).}
A seleção por roleta requer valores positivos proporcionais à qualidade a serem
maximizados. Como o RMSE é uma métrica de custo (minimização), converte-se o
erro em uma aptidão estritamente decrescente $\phi_i = 1 / (1 + \mathrm{RMSE}_i)$.
Dessa forma, indivíduos com menor erro recebem fatias proporcionalmente maiores
na roleta (por exemplo, um indivíduo perfeito com $\mathrm{RMSE}=0$ obtém
$\phi=1.0$, enquanto um com $\mathrm{RMSE}=99$ recebe apenas $\phi=0.01$). A
probabilidade de seleção do indivíduo $i$ em uma população de tamanho $P$ é:
\begin{equation}
  p_i = \frac{\phi_i}{\sum_{j=1}^{P} \phi_j}.
\end{equation}
Quando a média dos erros populacionais é elevada, as diferenças relativas
entre $1/(1+\text{RMSE})$ tornam-se insignificantes, reduzindo a pressão
seletiva.

\paragraph{2. Seleção por Torneio.}
Amostram-se aleatoriamente $k$ indivíduos com reposição e seleciona-se aquele
com o menor RMSE. O parâmetro $k$ regula diretamente a intensidade da pressão
seletiva ($k=2$ é o padrão e $k=3$ impõe maior rigor).

\paragraph{3. Seleção $\eps$-Lexicase (MAD).}
Diferente dos métodos escalares, a seleção lexicase~\cite{helmuth2015lexicase}
trabalha diretamente com a matriz de erros individuais $E \in \mathbb{R}^{P \times N}$,
em que $e_t(i) = |\mathrm{Eval}(\mathit{Ind}_i, \mathbf{x}_t) - y_t|$. O algoritmo
executa os seguintes passos para selecionar cada progenitor:
\begin{enumerate}
  \item O conjunto de candidatos sobreviventes inicia com toda a população
        ($C = \{1, \dots, P\}$) e os $N$ casos de treino são embaralhados
        aleatoriamente.
  \item Para o caso $t$, calcula-se o menor erro entre os candidatos atuais:
        $e^*_t = \min_{i \in C} e_t(i)$.
  \item Na variante $\eps$-lexicase para regressão~\cite{lacava2016epsilon},
        sobrevivem ao caso $t$ todos os indivíduos cujo erro satisfaça:
        \begin{equation}
          e_t(i) \le e^*_t + \lambda(e_t),
          \label{eq:eps_lexicase}
        \end{equation}
        onde o limiar adaptativo $\lambda(e_t)$ é o Desvio Absoluto Mediano
        (MAD) dos erros dos candidatos no caso $t$:
        \begin{equation}
          \lambda(e_t) = \operatorname{mediana}_{j \in C}\Big(\big|e_t(j) - \operatorname{mediana}_{k \in C}(e_t(k))\big|\Big).
        \end{equation}
  \item O processo itera sequencialmente pelos casos. Se $|C| = 1$, o indivíduo
        remanescente é retornado imediatamente. Se os $N$ casos forem
        esgotados com múltiplos sobreviventes, sorteia-se um deles
        uniformemente.
\end{enumerate}
Esse mecanismo preserva linhagens que se destacam em subconjuntos distintos
de dados, evitando que soluções com erro médio medíocre, mas cruciais em
casos raros, sejam prematuramente eliminadas.

\paragraph{Elitismo e Desempate Parcimonioso.}
O arcabouço implementa elitismo (\code{elite\_size=1}) preservando o melhor
indivíduo absoluto de treino. Em caso de empate de aptidão dentro de uma
tolerância relativa de $10^{-6}$, o desempate prioriza a árvore com menor
número total de nós.

\paragraph{Pressão de Parcimônia Lexicográfica (\emph{Bloat Control}).}
Inspirado em \citet{luke2002lexicographic}, implementou-se um desempate
anti-inchaço na seleção: quando a $\eps$-lexicase atinge $\le 4$ candidatos
finais após avaliar ao menos $25\%$ dos casos de treino, escolhe-se o
indivíduo com menor tamanho em nós dentre aqueles cuja aptidão global esteja
dentro de $5\%$ do melhor candidato.

\paragraph{Rejeição de Duplicatas (\code{dedup\_offspring}).}
Para impedir que clones inertes ocupem vagas na nova geração, cada descendente
tem sua forma canônica confrontada com as já aceitas na geração atual; caso já
exista indivíduo idêntico, o descendente sofre mutação obrigatória (até 3
tentativas) antes da inserção.

% -----------------------------------------------------------------------------
\subsection{Coleta Estatística e Verificação de Software}
\label{sec:impl_testes}

Em todas as gerações, o módulo \code{stats.py} registra: erro mínimo, médio,
máximo e desvio-padrão; contagem de indivíduos únicos e repetidos; número de
filhos de crossover melhores e piores que a média e que o melhor pai; eficácia
das mutações; e tamanho médio, máximo e profundidade média da população.